\section{Conclusion}
\label{sec:conclusion}

Spectral compression can preserve the likelihood information needed for nearly optimal community recovery. In a full-rank SBM, the membership subspace contains the block-constant likelihood contrasts. First-order analysis of the empirical adjacency eigenpairs makes this relation quantitative at logarithmic density, and an admissible fixed positive floor preserves the leading Chernoff exponent. The subsequent likelihood and EM analysis controls the dependent whole-graph profiles.

For $K=2$, the original growing likelihood procedure has an end-to-end guarantee. For general fixed $K$, Algorithm~1 supplements the growing fit with $\lceil\log n\rceil$ independent likelihood-residual++ initializations, compares the fitted candidates using the same observed working mixture likelihood, and performs a final EM update before decoding. The resulting recovery bound is
\[
 \mathbb E\,\Mis(\widehat z,z)\le e^{-(1+o(1))I_n}.
\]
Initialization and likelihood selection are part of the proof; weak recovery is not an assumption imposed on an unspecified initializer. The guarantee matches the leading oracle exponent and yields exact recovery above the corresponding logarithmic information threshold.

The paired 240-graph experiments evaluate the original growing and parameter-replacement procedures. They support the usefulness of all-node likelihood gains and likelihood-accepted repairs in resolving severe initialization failures. A full evaluation of the safeguarded combination remains to be carried out. Removing the safeguard while retaining the general fixed-$K$ theorem, establishing calibrated spectral concentration for a lower-density regularizer, and reducing the cost of full profile reconstruction are separate extensions of the completed logarithmic-density theory.
